\subsection{3. Die neue präzise
Händigkeitsschranke}\label{die-neue-pruxe4zise-huxe4ndigkeitsschranke}

Die eingefrorene Tabelle in \codeword{field\_dictionary.py}, Zeile 327,
bezeichnet den zweiten Faktor als
\codeword{f\textasciicircum{}dagger (1/2,0)}. Dafür muss die
Feldabbildung explizit festgelegt werden. Für ein tatsächliches
adjungiertes linkes Weylfeld gilt dagegen ψ† in (0,1/2). Hermitesche
Konjugation vertauscht diese Darstellungen; siehe
\href{https://arxiv.org/pdf/0812.1594\#page=9}{Dreiner, Haber und
Martin, Abschnitt 2}.

In diesem \textbf{wörtlichen Feldadjungierten-Fall} ist das Produkt

\[\adjustbox{max width=\linewidth}{$\displaystyle 
 (1,0)\otimes(0,1/2)=(1,1/2)
$}\]

ein irreduzibler Lorentzraum mit komplexer Dimension sechs. Unter der
Rotationsuntergruppe zerfällt er weiterhin in Spin 3/2 und Spin 1/2;
diese beiden Rotationsräume sind aber nicht getrennt boostinvariant. Der
Prüfer verifiziert den skalaren Kommutanten sämtlicher sechs Generatoren
des komplexifizierten Lorentzraums; die beiden unabhängigen
sl(2)-Faktoren wirken als irreduzibles äußeres Tensorprodukt.

Ein konkreter Zeuge genügt bereits. In obiger Basis setze

\[\adjustbox{max width=\linewidth}{$\displaystyle 
 J_z=J_z^{(1)}\otimes I_2+I_3\otimes J_z^{(1/2)},\qquad
 K_z=J_z^{(1)}\otimes I_2-I_3\otimes J_z^{(1/2)}.
$}\]

Bis auf den konventionellen Faktor i ist K\_z der Boostgenerator für
entgegengesetzte Händigkeit. Dann gilt exakt

\[\adjustbox{max width=\linewidth}{$\displaystyle 
 [P_{1/2},J_z]=0,\qquad
 [P_{1/2},K_z]=
 \begin{pmatrix}
 0&0&0&0&0&0\\
 0&0&4/3&0&0&0\\
 0&-2/3&0&0&0&0\\
 0&0&0&0&2/3&0\\
 0&0&0&-4/3&0&0\\
 0&0&0&0&0&0
 \end{pmatrix}.
$}\]

Dieser Kommutator hat Rang vier. Der Rotationsprojektor kann deshalb in
diesem Vertrag nicht als Lorentz-kovarianter Weyl-Ausleseprojektor
verwendet werden.

\textbf{Das ist kein pauschaler Ausschluss des ursprünglichen
Modenkomposits.} Ein bloßer Fock-Erzeuger f† ist noch kein adjungiertes
lokales Weylfeld. Bereits die übliche linke Weylfeldentwicklung enthält
Erzeugungs- und Vernichtungsterme mit passenden Spinorwellenfunktionen;
siehe \href{https://arxiv.org/pdf/0812.1594\#page=24}{Dreiner, Haber und
Martin, Gleichungen 3.1.3-3.1.5}. Eine unabhängige oder passend
ladungskonjugierte gleichhändige Komponente kann den gleichhändigen
Projektor erhalten. Sie muss aber als konkreter Adapter mit korrekter
Ladung, CAR, W-Kontraktion und Dynamik gezeigt werden. Aus dem Symbol †
allein folgt weder ihr Vorliegen noch ihr Ausschluss.

\subsection{4. Konkreter nächster
Annahmetest}\label{konkreter-nuxe4chster-annahmetest}

Der Anschlussauftrag ist jetzt endlich und entscheidbar formuliert:

\begin{enumerate}
\def\labelenumi{\arabic{enumi}.}
\tightlist
\item
  Den tatsächlichen Feldadapter für \codeword{D\_r ∼ b f†}
  beziehungsweise den davon verschiedenen \codeword{χ† ∼ b†ηf†}
  hinschreiben: gepunktete/ungepunktete Indizes, Ladung, Fourieranteile,
  erlaubte Zusatzkomponenten.
\item
  Für genau diesen Adapter alle Rotations- \textbf{und Boostgeneratoren}
  auf dem zusammengesetzten Träger berechnen. Ein Kandidat mit
  Lorentz-Auslese C muss \codeword{C L\_composite = L\_Weyl C} erfüllen.
  Der neue gleiche-/gegengängige Test liefert dafür positive und
  negative Referenzfälle.
\item
  Die Kontraktion am unveränderten W ausführen und anschließend das
  tatsächliche Matrixelement \codeword{C D\_r†|Ω>} beziehungsweise den
  entsprechenden Ladungskanal berechnen. Erst dessen normiertes
  Spektralgewicht kann den vorhandenen geladenen Pol mit einem
  Spin-1/2-Feld verbinden.
\item
  Falls der wörtlich adjungierte Fall gewählt wird, keinen bloßen
  Rotationsprojektor als Lorentzprojektor verwenden. Ein zusätzlicher
  Impuls-/Ableitungsadapter oder eine andere Händigkeit kann untersucht
  werden, zählt aber als ausdrückliche neue Struktur.
\end{enumerate}

Dieser Test ergänzt die gemeinsame Schattenrekonstruktion: Die dort
bewiesene vollständige innere Tomographie ersetzt die hier fehlende
Lorentz-Intertwining-Abbildung nicht. Umgekehrt definiert ein richtiger
Lorentzprojektor noch kein ausführbares natives Messinstrument.

\subsection{Reproduktion}\label{reproduktion}

\begin{verbatim}
python3 -B replay_late_field.py
\end{verbatim}

Benötigt werden nur Python und SymPy. Der separate Lauf verifiziert zwei
Quellenpins sowie \textbf{36 exakte Bedingungen pro Modus}, keine
numerischen Bedingungen. Normaler Lauf und \codeword{-OO} ergeben
byte-identische JSON-Ausgaben; Warnungen gelten als Fehler und alle
Guards bleiben unter Optimierung aktiv. Die vollständigen Matrizen
stehen in \codeword{late\_field.normal.json}.

Das Ergebnis-JSON hat SHA-256
\codeword{cfd9e15cec5c2b0ce3de9397614836b16b24b4adb8f6711eb704925589965774}.
Der Prüfer hat SHA-256
\codeword{f18e6fd2af7c1a28a81e85654df68fe805e59b1d957801611d2668678550dfc2}.
Die fremde Feldsuite wird nur eingefroren und punktuell unabhängig
geprüft, nicht als vollständig fehlerfreie Gesamtreproduktion
ausgegeben.

\begin{center}\rule{0.5\linewidth}{0.5pt}\end{center}

\section{Nachtragsprüfung C: fundamentale Reduktion und globale
Quelle}\label{nachtragspruxfcfung-c-fundamentale-reduktion-und-globale-quelle}

\section{Nachtrag C: überwiegend bereits abgedeckte fundamentale
Reduktion}\label{nachtrag-c-uxfcberwiegend-bereits-abgedeckte-fundamentale-reduktion}

\subsection{Übernahmeurteil}\label{uxfcbernahmeurteil}

Der neue Text ist eine ausführlichere Fassung der bereits eingefrorenen
fundamentalen Reduktion, ergänzt um die weitgehend gleiche frühere
Gesprächszusammenfassung. Er bringt \textbf{keinen neuen H-Term, keine
neue Operationsquelle und keinen neuen Transferbeweis}. Seine
wesentlichen mathematischen Einschränkungen sind richtig und bleiben in
der gemeinsamen Darstellung erhalten. Eine neue Großrechnung ist dafür
nicht notwendig.

Neue Eingabe: 299 Zeilen, SHA-256
\codeword{a10f702306c7409535df9f4d26f0995970833a62297a1cb93d8c95cdce746d4d}.
Quelle:
\codeword{/Users/stefanhamann/.codex/attachments/b0d7e533-3ca8-4a84-966c-2546b066dd08/pasted-text.txt}.

Bereits eingefrorene Kurzfassung:
\codeword{universalraum-singlet-observable-20260915/late\_sources/fundamental\_reduction.txt},
99 Zeilen, SHA-256
\codeword{0081cebd9759e8fa776995ee8efb3231c7611a82bea9cdeea86967073d1459a2}.
Die Dateien sind nicht bytegleich; der Neuheitsbefund betrifft ihren
Inhalt.

\subsection{Was bestätigt und bereits abgedeckt
ist}\label{was-bestuxe4tigt-und-bereits-abgedeckt-ist}

{\def\LTcaptype{none} % do not increment counter
\begin{longtable}[]{@{}
  >{\raggedright\arraybackslash}p{(\linewidth - 2\tabcolsep) * \real{0.5000}}
  >{\raggedright\arraybackslash}p{(\linewidth - 2\tabcolsep) * \real{0.5000}}@{}}
\toprule\noalign{}
\begin{minipage}[b]{\linewidth}\raggedright
Aussage
\end{minipage} & \begin{minipage}[b]{\linewidth}\raggedright
Urteil und Grenze
\end{minipage} \\
\midrule\noalign{}
\endhead
\bottomrule\noalign{}
\endlastfoot
Volle Schur-Antwort T†F(H−E0)T=γF I64 & Richtig bei gruppeninvariantem H
und Grundzustand sowie irreduziblem ursprünglichem Entnahmemultiplett.
Gilt für die volle beschränkte Spektralfunktion, nicht nur einen
dominanten Pol. Schon in v1.6.7 BIG\_PICTURE B2 und der eingefrorenen
Kurzfassung enthalten. \\
Aus C\_rs(t)=δ\_rs c(t) folgt kein Ortsindex in diesen 64 Labels &
Richtig. Eine andere lineare Benennung erzeugt nur einen Gramfaktor.
Keine pauschale Sperre gegen Vielteilchendynamik, nichtinvariante
Referenzzustände oder andere Operatorfamilien. \\
Symmetriekommutant und Operationskommutant haben verschiedene Bedeutung
& Richtig und bereits unterschieden. Erstgenannter erlaubt Dynamik
zwischen Vorkommen desselben Typs; ein großer Kommutant eines
eingeschränkten Operationssatzes kann fehlenden Zugriff anzeigen. Seine
Größe beweist weder Raum noch ausführbare Kontrolle. \\
μN ist bei festem Gesamt-N eine Konstante & Richtig. Auf demselben
präparierten Eingang bleiben freie Heisenbergentwicklungen neutraler
Observablen identisch. Für ganze Interventionsfolgen muss auch deren
zugelassener Gesamtvertrag respektiert werden. \\
Eine globale Zustandsregel braucht keinen äußeren Präparator & Richtig.
Das Verbot, mit N-erhaltenden Operationen aus N=0 nach N=64 zu gelangen,
ist ein Satz über einen bestimmten Präparationsweg. Es verbietet weder
eine globale Randbedingung noch eine andere begründete Zustandsregel.
Interne Detektoren und Präparationen bleiben trotzdem herzuleiten. \\
Keine nichttriviale Fermion-Teilparität in einer nativen Bank bei festen
Bosonen & Bereits exakt geprüft: zusammenhängender W-Paargraph und
binärer Rang 63. Ein Paritätsverbot getrennt kopierter Banken darf nicht
auf beliebige Teilmengen derselben Bank übertragen werden. \\
\end{longtable}
}

Das eigene v1.6.8-Ergebnis ergänzt diese Diagnose: Der bedingte lokale
Phasenimpuls und die terminale Tomografie benutzen einen gemeinsamen
invarianten N=2-Dreizustandsraum. Dessen Startzustand ist nicht der in
Schurs Aussage vorausgesetzte invariante N=64-Grundzustand. Es gibt
daher keinen Widerspruch und keine Berechtigung, die neue Umverteilung
auf dessen geladenen Einlochpol oder auf räumliche Bewegung zu
übertragen.

\subsection{μ=Δ/50: gegen die älteren gepinnten Zertifikate
bestätigt}\label{ux3bcux3b450-gegen-die-uxe4lteren-gepinnten-zertifikate-bestuxe4tigt}

Das ältere
\codeword{universalraum-native-ground-response-20260915/ground\_replay/weak\_coupling\_ground\_normal.json}
enthält genau die genannte Grenze: g/Δ=1/20, μ/Δ=1/50 und neuer
Grundsektor N=0. Sein SHA-256 ist
\codeword{e1cd19988edd33c6799ec6b80f0a52d56e1a532bfe1b56870623a4efaf06bd0c};
dies stimmt mit \codeword{ground\_replay\_manifest.json} überein.
Normaler und optimierter vorhandener Bericht sind weiterhin bytegleich.
Der eingetragene Programmpin
\codeword{794554393c495ba395f34e6a58d490da821408d4f8effbc394653cd4cec4e80a}
stimmt mit der vorhandenen
\codeword{ground\_replay/work/many\_pair/weak\_coupling.py} überein.

Die stärkere und direktere Begründung ist bereits in
\codeword{universalraum-v16-integrated-20260915/RESULTS.md} §3 und
dessen \codeword{new\_checks.json} dokumentiert. Der Checkerhash
\codeword{c3304a1cf502e5c4fffd63391cfb29db02ad3f2b3517f96e047c07170b63416e}
stimmt mit dem vorhandenen Programm und \codeword{replay\_manifest.json}
überein. Aus der dort geprüften Voll-Fock-Schranke

\[\adjustbox{max width=\linewidth}{$\displaystyle 
A=\sum_A P_A^\dagger P_A\le\frac{15}{2}N_f
$}\]

und quadratischer Ergänzung folgt

\[\adjustbox{max width=\linewidth}{$\displaystyle 
H_\mu\ge\left(\mu-\frac{15g^2}{2\Delta}\right)N_f+2\mu N_b.
$}\]

Die einzige erneut benötigte Rechnung ist rational:

\[\adjustbox{max width=\linewidth}{$\displaystyle 
\frac1{50}-\frac{15}{2}\left(\frac1{20}\right)^2=\frac1{800},
\qquad2\mu/\Delta=\frac1{25}.
$}\]

Also gilt \(H_\mu\ge\Delta N/800\). Da N=0 nur das leere Fockvakuum
enthält und Hμ dieses mit Energie null annihiliert, ist es eindeutig
minimal. Der globale Zustandsgegenvergleich ist damit bestätigt. Der
vollständige ältere Grundzustands-/Casimir-Prüflauf wurde für diesen
Nachtrag nicht neu ausgeführt; Pins, vorhandene Aussagen und die
relevante exakte Ungleichung wurden geprüft.

\subsection{Zwei Präzisierungen bei der
Übernahme}\label{zwei-pruxe4zisierungen-bei-der-uxfcbernahme}

\begin{enumerate}
\def\labelenumi{\arabic{enumi}.}
\tightlist
\item
  In §7 sollte „ein positiver Unterschied`` durch \textbf{„eine
  nichtverschwindende Differenz``} ersetzt werden. Ein negativer
  signierter Unterschied ist ebenso ein kausaler Interventionsnachweis;
  genau dies zeigt der neue N=2-Phasenversuch.
\item
  Die globale Sektorwahl und die Beobachtbarkeit innerhalb eines festen
  Gesamtsektors getrennt halten. Ein System-Referenz-Vergleich kann
  relative Ladungsenergien aufdecken; ein Zusatz μN\_total auf einem
  festgehaltenen Gesamtsektor bleibt dennoch eine unbeobachtbare
  gemeinsame Phase.
\end{enumerate}

\textbf{Fazit:} Hauptpunkt abgedeckt, keine neue globale Schließung. Die
ausführlichere Quelle stärkt die Dokumentation der Voraussetzungen und
die Abgrenzung der Zustandsregel. Die tatsächliche neue Konstruktion
dieser Runde bleibt die gemeinsame bedingte N=2-Ausführung aus kausalem
Eingriff und terminaler Zustandstomografie.

\begin{center}\rule{0.5\linewidth}{0.5pt}\end{center}

\section{Historischer Teil: vollständige Hauptfassung
v1.6.7}\label{historischer-teil-vollstuxe4ndige-hauptfassung-v1.6.7}

Die folgende Fassung bleibt wortgetreu erhalten. Den aktuellen Stand
liefern die vorangestellten Kapitel v1.6.8. Frühere Aussagen, Fragen und
Vorhaben werden dadurch nicht rückwirkend neu bewiesen. Die neue
Synthese ist separat vollständig im eingefrorenen Prüfpaket enthalten.

\section{TFPT / Universalraum: die richtige einfache
Kette}\label{tfpt-universalraum-die-richtige-einfache-kette}

\subsection{Vollständige Konsolidierung v1.6.7 - 15. September
2026}\label{vollstuxe4ndige-konsolidierung-v1.6.7---15.-september-2026}

\subsubsection{Ergebnis und
Leseschlüssel}\label{ergebnis-und-leseschluxfcssel-1}

Die zuletzt ergänzte Vogelperspektive steht im Teil B dieser Fassung:
\textbf{Gesucht wird der kleinste gemeinsame beobachtbare Prozess hinter
den verschiedenen Ansichten.} Ein nativer Rang-60-Kodierer ist bereits
vorhanden. Seine Informationsverluste, die minimalen
symmetrieverträglichen Dynamiken und der Unterschied zwischen
Korrelation und kausalem Eingriff wurden neu berechnet. Das verschiebt
die fundamentale Priorität von weiteren lokalen Näherungswerten zur
Herkunft der Operationen und zur gemeinsamen Schattenkarte.

Die neue Runde schließt zwei kleine Strukturfragen vollständig: Der
symmetrische Zustandsraum bei zwei Bosonen hat exakt vier Dimensionen;
der bisher unbestimmte Rest von 4035 bilinearen Komponenten zerfällt
vollständig in drei Darstellungen. Die vier Zustandsrichtungen und ihre
Kopplungen wurden zusätzlich am ursprünglichen Tensor berechnet. Daraus
folgt ein vollständiger sechs-dimensionaler Ausschnitt.

Die gewünschte einfache Kette existiert ebenfalls, aber anders als
vorgeschlagen: Man muss den \textbf{ganzen Hamiltonoperator}
wiederholen, nicht Bosonzahl und Iterationsnummer gleichsetzen. Die
ersten fünf Lanczos-Glieder sind jetzt exakt bekannt. Sie bestimmen zehn
Energiemomente des gefüllten Startzustands und verschärfen die obere
Grundzustandsgrenze geringfügig.

\textbf{Das ist keine vollständige TOE.} Insbesondere fehlen die
physische Auswahl des Modellvertrags, ausführbare ursprüngliche
Instrumente, derselbe geladene Transport im gemeinsamen Träger und die
relativistische Kontinuumskonstruktion. Dass die bisherigen Grunddaten
dafür noch nicht eindeutig genügen, wird unten an einer einfachen
Familie verschiedener Dynamiken nachgewiesen.

Kennzeichnungen: \textbf{exakt} bezeichnet eine bewiesene Aussage im
genannten Modell; \textbf{bedingt} eine Aussage mit zusätzlich gewährter
Operation oder bekanntem früherem Satz; \textbf{numerisch} eine
Näherung; \textbf{offen} einen fehlenden Nachweis. Komplexe
Algebradimensionen zählen weder physikalische Raumrichtungen noch
unabhängig hergeleitete Teilchen oder ausführbare Operationen.

Die Fassung baut auf dem vollständigen Hauptdokument v1.6.6 auf. Die
aktuelle Gesamtdarstellung steht voran; die frühere Fassung
einschließlich ihrer Herleitungen, Gegenbeispiele und älteren Anhänge
bleibt vollständig erhalten. Alle sechs in dieser Runde eingesandten
Texte werden unverändert dokumentiert und hier mit überprüften
Ergänzungen und ausdrücklich benannten Korrekturen gelesen.

